\begin{textsample}{POS Dim 1 – pb – Score 38.00 – pb/pb\_em\_6.txt}  \label{ex:f1_pos_130}
1.3 PRÁTICAS EDUCATIVAS DA PROPOSTA CURRICULAR DO ENSINO MÉDIO DA PARAÍBA 1.3.1 PROCESSOS \textbf{METODOLÓGICOS} Na Educação os resultados não são alcançados de \textbf{uma} hora para outra, há \textbf{um} tempo para o aprendizado.
Tem início, meio e fim.
Mas pensemos em “ fim ” não como término, mas como finalidade, ou meta, pois quando há algo a ser pretendido, alcançado, é essencial construir o projeto para se chegar no objetivo, planejando as maneiras de como caminhar, antes mesmo de iniciar os primeiros passos.
Sabemos, também, \textbf{que} com a vivência da caminhada, o próprio terreno no qual se caminha também pode se manifestar diferente, \textbf{exigindo}, portanto, diferentes formas de continuar caminhando.
Entretanto, como trilhar essa caminhada é o \textbf{método} \textbf{que} pode ser utilizado ao longo do percurso.
Assim, o caminho é a Educação, a meta é a aprendizagem, os veículos para trilhar a caminhada são os \textbf{métodos}.
Em toda a história da Educação no Brasil passamos por \textbf{inúmeros} meios ( metodologias ) de ensino até os mais recentes estudos e práticas educacionais da atualidade.
E, ultimamente, com as emergentes necessidades \textbf{imediatas} de aprendizagem pelas novas tecnologias educacionais, devido às grandes transformações do ensino presencial e remoto em tempos de pandemia, professores/as, alunos/as, gestores/as e comunidade escolar em geral, em todas as \textbf{instâncias}, estão passando por profundas transformações com resultados ainda indefinidos.
Nesse sentido, \textbf{vivemos} desastrosos e assustadores tempos de catástrofe pandêmica em 2020, mas sabemos também \textbf{que} a Educação não pode parar.
Assim, devido às determinações da lei para o isolamento e o distanciamento social, as escolas ficaram com suas portas, literalmente, fechadas e as aulas, em muitas escolas, tiveram \textbf{que} ser transmitidas remotamente.
Ainda, muitos professores buscam maneiras para aprender a manipular aplicativos de smartphones, softwares e conhecer plataformas e estratégias para aulas síncronas e assíncronas em AVA ( Ambientes Virtuais de Aprendizagem ).
É importante \textbf{lembrar} também \textbf{que} a escola é \textbf{uma} parte da Educação, visto \textbf{que} não faz educação sozinha e nem a educação pode ter início na escola.
Ao nascer, o indivíduo já inicia a sua aprendizagem na escola da vida, muito embora as pessoas deleguem à escola formal, ao colégio, todas as atribuições da educação do ser humano.
Nesse parâmetro, a escola \textbf{precisa} de integração por parte das famílias dos/as estudantes para \textbf{que} possam somar e afinar a educação formal com a educação familiar, visando o bem comum com vistas ao aprendizado conectado com a vida, sem a barreira entre \textbf{teoria} ( estudo formal ) e prática ( o dia a dia do indivíduo ).
Na realidade, a escola deve oferecer as vivências \textbf{teóricas} e práticas harmonicamente conectadas para \textbf{que} o estudante \textbf{veja} sentido no \textbf{que} é estudado, concebendo os saberes apreendidos nesse espaço como caminhos e apontamentos para toda a vida.
Se isso fica claro para eles, o processo ensino-aprendizagem se concretiza no estudante.
A Paraíba, em seus 223 municípios, apresenta realidades diversas em suas localidades e isso precisa ser respeitado.
São zonas rurais, urbanas, comunidades indígenas, quilombolas, ciganas, praieiras, \textbf{ribeirinhas} e não há \textbf{um} parâmetro único para delimitar \textbf{uma} metodologia de ensino \textbf{que} abarque tantas realidades.
Para José Carlos Libâneo ( 1994, p.178 ) : “ aula é toda situação didática na qual se põem objetivos, conhecimentos, problemas, desafios com fins instrutivos e formativos, \textbf{que} incitam as crianças e jovens a aprender".
Assim, entende -se \textbf{que} cada aula é especialmente única, com suas próprias características e deve ser pensada e planejada de acordo com as realidades e juventudes, suas respectivas diversidades e necessidades observadas nos/as estudantes \textbf{que} irão recebê -la.
Sendo assim, o estudo das características peculiares de cada localidade é \textbf{que} vai ser a linha mestra para determinar especificamente os meios para atingir os resultados da aprendizagem desejada, muito embora \textbf{possamos} traçar alguns caminhos para encontrar os procedimentos \textbf{metodológicos} mais adequados a cada realidade.
Nessa perspectiva, a metodologia de ensino se faz \textbf{um} conjunto de princípios e/ou diretrizes sociopolíticos, epistemológicos e psicopedagógicos capaz de encaminhar procedimentos orgânicos e sequenciados para orientar o processo ensino-aprendizagem e direcionar o estudante paraibano aos objetivos almejados.
Pretendemos \textbf{que} os estudantes da Paraíba sejam proativos e criativos, portanto, as escolas precisam adotar metodologias em \textbf{que} eles se envolvam em atividades cada vez mais complexas para \textbf{que} possam tomar decisões, avaliar resultados e experimentem \textbf{inúmeras} possibilidades de mostrar sua iniciativa.
Eles precisam encontrar na sua escola, aulas \textbf{que} os convidem a aprender ativamente com problemas reais, \textbf{que} façam parte de suas vidas, de seus cotidianos.
Quando o professor trabalha com jogos, desafios, atividades variadas, leituras estimulantes para os estudantes e outras práticas, ele supera aquela aula “ presa ”, maçante, engessada pelos conteúdos do livro didático.
Precisamos abolir o antigo professor \textbf{que} luta para “ vencer o conteúdo ”.
E o termo “ abolir ” cai muito bem \textbf{aqui}, pois é \textbf{uma} abolição da escravatura \textbf{que} \textbf{aqui} se propõe.
É \textbf{preciso} libertar -se de \textbf{um} modelo único.
Ou existe \textbf{um} único livro didático \textbf{que} atinge toda a pluralidade do Estado?
Não confundamos o \textbf{que} se propõe \textbf{aqui}.
O livro não deve ser extinguido.
O livro continua sendo \textbf{um} dos caminhos para a educação, mas outras leituras devem ser descobertas, estimuladas e respeitadas, como a tradição oral de \textbf{uma} comunidade, a leitura simbólica e imagética, a leitura das diferentes concepções de tempo e espaço, individual e coletivo, entre outras.
O professor não pode mais pensar em “ vencer \textbf{um} conteúdo ”. \textbf{Educar} não é brigar com conteúdos.
O desafio não é esse.
Os objetos de conhecimento ( conteúdos ) devem ser saboreados pelos/as estudantes, e, quando isso ocorre, a experiência leva esse sabor.
Aprender é prazeroso tanto quanto ensinar.
Mesmo porque o ato de ensinar carrega dentro de si o ato de aprender.
Quem ensina também está aprendendo, assim como quem está aprendendo está ensinando. \textbf{Educar} não é \textbf{uma} \textbf{mera} transferência de conceitos, informações, experiências e conteúdos didáticos.
Como afirma Paulo Freire na sua \textbf{Pedagogia} da autonomia ( 2007, p. 22 ) : “ [... ] \textbf{educar} não é transferir conhecimento, mas criar possibilidades para a sua própria produção ou a sua construção ” Essa visão de acúmulos externos \textbf{que} devem ser depositados nas \textbf{mentes} dos estudantes não passa de \textbf{uma} ilusão.
O desenvolvimento de \textbf{um} ser não deve ser mensurado pela quantificação de informações, mas pela qualidade real de seu aprendizado.
Os valores quanto ao \textbf{que} se ensina e como se ensina precisam ser constantemente questionados pelo profissional da educação.
Assim como os conteúdos são propostos a partir das habilidades e competências a serem alcançadas, a forma, ou melhor, as formas com \textbf{que} serão trabalhadas também devem ser planejadas neste mesmo trilho.
Primeiro pensa -se o quê e o porquê \textbf{uma} meta deve ser atingida, em seguida, como dever ser trilhada.
Esse como são os processos \textbf{metodológicos}.
A implantação de metodologias inovadoras requer \textbf{uma} gama de transformações, mas a mudança inicial é de mentalidade.
Na verdade, \textbf{necessita} -se de \textbf{uma} transmutação, trans ( de \textbf{um} \textbf{lado} para outro ), mutare ( mudar ), \textbf{uma} vez \textbf{que} é \textbf{uma} mudança de dentro para fora.
Cada indivíduo, mesmo \textbf{que} estimulado por condicionantes externos, só pode fazer essa alteração por ele próprio.
Mudar a forma de pensar, falar, agir, sentir e fazer, em se tratando de processos \textbf{metodológicos}, pode ter seu primeiro passo no adequar -se aos novos tempos.
Em geral, a imensa maioria dos profissionais da educação no Brasil não acompanhou o processo natural de evolução.
Há \textbf{uma} dissonância entre professores e estudantes.
Os professores \textbf{que} ainda ensinam do mesmo modo de décadas, e até \textbf{séculos} atrás, estão no passado. \textbf{Enquanto} \textbf{que} os seus estudantes, \textbf{que} estão no presente, não suportam a ideia de \textbf{que} alguém \textbf{que} \textbf{vive} no passado possa instruí -lo para o futuro.
Esse abismo é, no mínimo, incoerente.
Aprender com o \textbf{que} passou traz experiência, mas experiente não quer \textbf{dizer} manter práticas e \textbf{métodos} \textbf{que} deram certo na sua época \textbf{enquanto} estudante, ou quando começou sua carreira, quer \textbf{dizer} \textbf{que} aprendeu com experiências, inclusive \textbf{que} pode continuar a aprender com novas experiências a todo momento. \textbf{Um} detalhe importante é \textbf{que} metodologias inovadoras estão a serviço não pura e simplesmente a partir das competências digitais com o uso de novas tecnologias, mas com \textbf{uma} série de práticas \textbf{metodológicas} \textbf{que} muitas vezes nem precisem utilizar energia elétrica.
Portanto, não dispor de equipamentos tecnológicos de última geração não impede a utilização de metodologias inovadoras.
Do mesmo modo, \textbf{uma} sala de aula convencional, com sua disposição usual de objetos e de pessoal, pode ser simplesmente alterada de acordo com a metodologia a ser empregada.
Para a utilização de novas metodologias, o professor e a equipe gestora da escola precisam estudar as possibilidades já existentes, como as Metodologias Ativas, por exemplo, para, a partir de então, poderem encontrar aquela, ou aquelas, \textbf{que} possam ser adequadas a cada fim.
Mas vale ressaltar \textbf{que} não existe \textbf{um} modelo pronto, exato, igual, imutável, para cada situação, entretanto, algumas virtudes são vitais para esses novos aprendizados : \textbf{vontade}, criatividade, bom senso, planejamento e a busca de sempre querer melhorar.
Assim, só é capaz de mudar aquele \textbf{que} se dispõe a tal.
Vide o \textbf{que} estamos passando em tempos de pandemia.
Aqueles \textbf{que} não se propõem a mudar e se readaptar \textbf{perdem} seu espaço, em qualquer \textbf{que} seja sua área.
E para isso é \textbf{preciso} muita \textbf{vontade}, inclusive, para quebrar resistências \textbf{que} impedem sua progressão, pois criatividade é \textbf{inerente} do ser humano.
Para o professor é \textbf{um} pulsar a cada aula planejada, a cada ideia concebida, a cada prática realizada, mas é \textbf{preciso} desenvolvê -la e, para alguns, ainda é necessário se descobrir \textbf{enquanto} ser criativo, pois um/a professor/a não é \textbf{um} \textbf{mero} repetidor.
Nessa perspectiva, como \textbf{veremos} a seguir, estamos constantemente nos autoavaliando e devemos fazê -lo a cada decisão usando de bom senso.
Através do uso de novas metodologias, o professor aprende a deixar de ser aquele \textbf{que} tende a explicar tudo, \textbf{exigindo} \textbf{que} o/a estudante anote tudo, pesquise e apresente o \textbf{que} memorizou.
O papel dele passa a ser aquele \textbf{que} media entre o conhecimento e o/a estudante, podendo estimulá -lo a ir além do \textbf{que} \textbf{conseguiu} fazê -lo sozinho.
Nessa \textbf{direção}, o professor precisa “ organizar as interações e atividades de modo \textbf{que} cada aluno se defronte constantemente com situações didáticas \textbf{que} lhe sejam as mais fecundas ” ( Perrenoud, 1995, p.28 ).
E, para isso, os estudantes precisam ser provocados.
Assim como \textbf{lembra} Rubem Alves “ \textbf{uma} concha só cria \textbf{uma} pérola porque esta foi provocada por \textbf{um} grão de areia ”.
Pérolas lindas podem brotar de cada estudante.
Portanto, a aprendizagem se constrói em processos \textbf{metodológicos} \textbf{que} precisam ser muito bem estudados e equilibrados no âmbito da aprendizagem personalizada, assim como da aprendizagem colaborativa, entre pares e por orientação, com profissionais mais experientes, como \textbf{um} especialista, \textbf{um} professor.
Para isso, os novos modelos de metodologias, como as Metodologias Ativas de Aprendizagem, trazem diversas possibilidades, como os jogos e as aulas roteirizadas com a linguagem de jogos ( gamificação ) ; aprendizagem ativa pela investigação ; aprendizagem \textbf{baseada} em problemas ; aprendizagem \textbf{baseada} em projetos dentro de cada componente curricular, projetos integradores ( interdisciplinares ), projetos transdisciplinares etc.
Todo professor da Paraíba precisa estudar algumas referências fundamentais como a \textbf{Pedagogia} da Autonomia, de Paulo Freire e a A aprendizagem significativa : a \textbf{teoria} de David Ausubel¹, por exemplo.
Sobre Metodologias Ativas de Aprendizagem, sugerimos o estudo da obra Metodologias Ativas para \textbf{uma} Educação Inovadora, dos \textbf{autores} Lilian Bacich e José Moran², \textbf{que} são referências no assunto.
Especificamente sobre as Metodologias Ativas no ensino presencial, híbrido e EAD o estudo do livro Games em educação : como os nativos digitais aprendem, de João Mattar³, pode ser \textbf{um} farol para esses novos tempos.
Não podemos deixar de mencionar \textbf{aqui} a importância, na qual todo este trabalho também tem se alicerçado e o professor precisa tê -lo sob seu domínio, \textbf{que} é o texto Educação : \textbf{um} tesouro a descobrir - relatório para a UNESCO, da Comissão Internacional sobre Educação para o Século XXI⁴.
Além desses modelos \textbf{que} o professor precisa estudar, há \textbf{uma} grande variedade de outros \textbf{que} podem servir de base para o desenvolvimento de novas metodologias.
Cabe ao professor mergulhar nesse universo e descobrir aquele(s) \textbf{que} melhor pode(m) alimentá -lo e conduzi -lo à meta principal, \textbf{que} é o aprendizado do estudante.
A Proposta Curricular do Ensino Médio da Paraíba apresenta \textbf{um} convite a você, leitor, para participar de \textbf{uma} construção de cidadãos para o mundo.
Construir \textbf{um} cidadão é incentivar a autonomia, é criar meios para \textbf{que} o/a estudante encontre a sua identidade na sociedade, para \textbf{que} ele seja \textbf{autor} ( e \textbf{ator} ) da sua própria história, a fim de não ser apenas “ mais \textbf{um} ” na multidão, mas “ \textbf{um} ” na multidão, não sendo coadjuvante, mas protagonista da sua história. ¹ Na obra David Ausubel, fica evidente a importância da valorização dos conhecimentos prévios dos/as estudantes, o \textbf{que} pode caracterizar \textbf{uma} aprendizagem prazerosa e eficaz ( AUSUBEL, D.
P.
A aprendizagem significativa : a \textbf{teoria} de David Ausubel.
São Paulo : Moraes, 1982 ). ² Os \textbf{autores} apresentam práticas pedagógicas \textbf{que} valorizam o protagonismo dos/as estudantes na construção do seus conhecimentos, assim como no desenvolvimento de competências essenciais para a educação nos tempos em \textbf{que} \textbf{vivemos} ( BACICH, Lilian ; MORAN, José.
Metodologias Ativas para \textbf{uma} Educação Inovadora : \textbf{uma} \textbf{abordagem} teórico-prática.
Porto Alegre : Penso, 2018 ). ³ Os \textbf{autores} apresentam práticas pedagógicas \textbf{que} valorizam o protagonismo dos/as estudantes na construção do seus conhecimentos, assim como no desenvolvimento de competências essenciais para a educação nos tempos em \textbf{que} \textbf{vivemos} ( BACICH, Lilian ; MORAN, José.
Metodologias Ativas para \textbf{uma} Educação Inovadora : \textbf{uma} \textbf{abordagem} teórico-prática.
Porto Alegre : Penso, 2018 ). ⁴ Essa obra, \textbf{que} conta com a contribuição de autoridades da educação do mundo todo, trata -se de \textbf{um} importantíssimo documento \textbf{que} cada professor/a deve se debruçar, pois traz os avanços tecnológicos do \textbf{século} XX e os princípios de coletividade no mundo \textbf{globalizado}, ressaltando a necessidade de trabalhar o ser humano em sua \textbf{totalidade}, numa visão cultural e multicultural. ( Educação : \textbf{um} tesouro a descobrir.
Relatório para a UNESCO, da Comissão Internacional sobre Educação para o Século XXI ).

% matched lemmas: abordagem, aqui, ator, autor, basear, conseguir, direção, dizer, educar, enquanto, exigir, globalizado, imediato, inerente, instância, inúmero, lado, lembrar, mente, mero, metodológico, método, necessitar, pedagogia, perder, possar, preciso, que, ribeirinho, século, teoria, teórico, totalidade, um, ver, vivar, viver, vontade
\end{textsample}
